\section{Theory, Guarantees, Failure Modes, and Reproducibility}
\label{sec:theory-repro}

\subsection{Guarantees are conditional statements}

A theoretical guarantee is meaningful only together with its assumptions. Smoothness, convexity, Lipschitz conditions, boundedness, stochastic assumptions, constraint qualifications, sampling assumptions, and access to derivatives or exact subproblem solutions determine what a theorem establishes. The atlas therefore does not reduce theory to a binary has-guarantee field. Instead, a theoretical statement is interpreted as a tuple
\begin{equation}
\mathcal{T}=\left(\mathcal{A},\mathcal{G},\mathcal{O}\right),
\end{equation}
where \(\mathcal{A}\) denotes assumptions, \(\mathcal{G}\) the guaranteed property, and \(\mathcal{O}\) the object to which the guarantee applies, such as iterates, expected regret, stationarity, convergence in probability, or approximation quality.

This distinction prevents guarantees from being compared outside their domains. A convergence result for a smooth convex objective, a global-convergence statement for a derivative-free polling framework, an asymptotic stochastic-search result, and a regret bound for online optimization answer different questions. None is converted into an unconditional empirical superiority claim.

\subsection{Theory and finite-budget behavior are complementary}

Asymptotic guarantees do not determine finite-budget performance, and strong finite-budget benchmark results do not establish a theorem. The evidence model therefore keeps theoretical and empirical axes separate. An algorithm may satisfy a convergence result while exhibiting unfavorable scaling under a particular budget, or may perform strongly on a benchmark despite having weaker available theory.

The distinction is especially important for stochastic and population-based optimization. Infinite-time convergence arguments, where available, need not predict behavior at \(100D\), \(300D\), or \(1000D\) evaluations. Conversely, a finite BBOB comparison cannot establish global convergence or refute a theorem whose assumptions and limiting regime differ from the experiment.

\subsection{Failure modes are mechanism--condition relations}

A failure mode is not defined as an algorithm being bad. It is a reproducible or documented relation between a mechanism and a condition under which an intended search behavior becomes ineffective, unstable, excessively costly, or insufficiently informative. We write this conceptually as
\begin{equation}
\mathcal{F}=\left(M,\theta,Y,P\right),
\end{equation}
where \(M\) is the implicated mechanism, \(\theta\) the condition, \(Y\) the observed outcome, and \(P\) its provenance. A mechanistic explanation requires more evidence than an observed deterioration.

The frozen literature matrix consequently uses bounded descriptions. High-dimensional simplex dynamics for Nelder--Mead, nonconvex curvature for BFGS, finite-budget cooling for simulated annealing, pheromone concentration for ant-colony optimization, and difficult Pareto-front geometries for decomposition or reference-direction methods are represented only at their supported evidence levels. Documented limitations such as known nonsmooth L-BFGS behavior are not generalized beyond the classes for which the evidence applies.

\subsection{From observation to mechanism claim}

The manuscript follows a claim-strength ladder. Descriptive algorithm identity occupies the lowest level and requires a canonical specification. Source-reported observations require the source and its conditions. Review synthesis requires comparable evidence across sources. A documented boundary requires a formal counterexample or adequate replicated empirical evidence. Project-validated conditional results require reproducible raw results, statistical analysis, and artifact provenance. A general mechanism claim requires cross-family evidence, tests of competing explanations, ablations, and external validation.

This ordering is important because performance deterioration alone cannot identify its cause. A higher-dimensional loss of target attainment could arise from search geometry, parameter scaling, population size, implementation overhead, termination, or the benchmark distribution. A mechanism explanation becomes credible only when alternatives are controlled or independently excluded. The present paper therefore does not promote its BBOB characterization to a general cross-family mechanism theory.

\subsection{Target semantics and authoritative measurement}

For the BBOB characterization, target attainment follows the frozen definition
\begin{equation}
f(x)\leq f_{\mathrm{opt}}+\Delta f,
\end{equation}
with the reporting grid \(\Delta f\in\{10^{2},10^{1},\ldots,10^{-8}\}\). The primary cross-function runtime unit is function evaluations divided by dimension, FE/D. Unreached targets remain unsuccessful observations rather than being removed from denominators.

The project deliberately separates measurement authorities. The COCO observer and logger are authoritative for benchmark target attainment and COCO-compatible postprocessing. The project objective recorder is authoritative for optimizer-independent objective-call accounting, best-so-far checkpoints, runtime metadata, and failures. The pinned public COCO Python interface does not expose a public \(f_{\mathrm{opt}}\) field suitable for a private reconstruction of exact target runtimes. Consequently, the analysis does not depend on undocumented internal optimum attributes or a hand-copied optimum table.

This split protects two different invariants: official benchmark semantics for target-runtime analysis and exact objective-call accounting for fixed-budget traces. Agreement between the recorder and COCO evaluation counts was checked in the publication executions.

\subsection{Failure-frontier discipline}

A failure frontier is useful only if its defining event is specified independently of the observations used to confirm it. Conceptually, for a performance requirement \(R\), a conditional frontier can be written
\begin{equation}
\partial\{\theta:R(A,\theta)=1\},
\end{equation}
where \(\theta\) indexes conditions such as dimension and budget. This notation does not by itself establish that such a frontier has been identified.

The project separated discovery instances 1--5 from held-out instances 6--10 and froze algorithms, dimensions, seeds, and FE/D budgets. However, the exact numerical reliability threshold and target subset required to convert target attainment into a binary frontier event were not frozen before the held-out campaign. Selecting those values afterward would turn a confirmatory boundary into a post hoc construction. The paper therefore reports the held-out campaign as project-validated characterization but does not claim an exact binary failure frontier.

This decision is a reproducibility result in its own right: experimental execution can be valid while a stronger interpretation remains unsupported because a decision rule was incompletely predeclared.

\subsection{Executable reproducibility}

Reproducibility is organized as a chain from scientific specification to executable evidence:
\begin{equation}
\text{protocol}\rightarrow
\text{configuration}\rightarrow
\text{execution freeze}\rightarrow
\text{raw artifact}\rightarrow
\text{analysis}\rightarrow
\text{claim}.
\end{equation}
Each link must remain inspectable. The repository therefore freezes benchmark software, comparator configurations, dimensions, instances, seeds, budgets, target semantics, statistical rules, and execution identifiers.

The production sequence also preserves historical corrections. Earlier execution freezes are not silently rewritten when a metadata or attribution defect is found. Instead, corrected freezes are versioned separately. This makes it possible to distinguish a scientific configuration change from a provenance or metadata correction.

For R10, the corrected cross-dimensional analysis contains 1,152 dimension--algorithm--function--budget cells. The EF004 production campaign produced 14,400 checkpoint rows and 64 algorithm-labelled COCO result roots across dimensions 5, 10, 20, and 40; objective-recorder and COCO counts agreed, and joint official \texttt{cocopp} postprocessing completed. The held-out EF005 campaign repeats the frozen design on disjoint instances for characterization.

\subsection{Evidence promotion and immutable freeze}

Not every executable output is publication evidence. Calibration, pilot, rehearsal, and run-complete-unverified states remain distinct from project-validated evidence. Failed, timed-out, and unobserved outcomes are retained rather than silently discarded. Statistical tables and figures are required to derive from promoted raw artifacts rather than hand-edited summaries.

Before manuscript drafting, the R15 audit checked execution-freeze integrity, source/fingerprint consistency, required interpretation records, and publication-control files. P031 then froze the principal evidence boundary using Git blob identities and validated that manifest in CI. Manuscript claims are consequently constrained by an immutable evidence snapshot.

The freeze does not imply that the atlas is complete. It means that the evidence supporting this version of the paper is fixed and auditable. A future experiment, corrected interpretation, or expanded literature audit must be introduced as a separately versioned amendment rather than silently changing the evidence beneath an existing claim.

\subsection{What reproducibility does not establish}

A reproducible experiment can still have limited external validity. Exact rerun capability does not prove that the chosen algorithms represent all optimization paradigms, that BBOB represents every application domain, or that a frozen parameterization is optimal for each method. Similarly, a canonical implementation does not guarantee fair computational cost across algorithms.

The paper therefore treats reproducibility as necessary for strong empirical claims but not sufficient for generality. Generalization depends additionally on benchmark coverage, mechanism diversity, parameter policy, cost accounting, independent replication, and consistency across conditions. This distinction motivates the bounded interpretation of the computational results in the next section.
